\section{Related Work}
\label{sec:related}

\paragraph{Multi-asset financial data generation.}
Generating realistic synthetic returns for multiple correlated assets has been approached from several angles. 
Classical methods include multivariate GARCH models \citep{engle2002dynamic, bollerslev1988capital}, copula-based generators \citep{marti2020corrgan, marti2021ccorrgan}, and factor bootstrap methods that combine rolling parameter estimates with resampled residuals. 
Among these, the factor bootstrap is surprisingly competitive: by decomposing returns into systematic (factor-driven) and idiosyncratic components, it preserves the cross-asset correlation structure without learning. 
Our work uses the factor bootstrap as the primary baseline and shows that a GAN can improve upon it---but only when the GAN's learning is carefully constrained.

\paragraph{GANs for financial time series.}
Several recent works apply GANs to financial data. 
MarketGAN \citep{huh2026marketgans} introduces a factor-structured generator with TCN backbone and WGAN-GP training, learning corrections to rolling OLS estimates of $\alpha$, $\beta$, and $\sigma$ simultaneously. 
Fin-GAN \citep{vuletic2023fingen} focuses on supervised probabilistic forecasting for individual series. 
Tail-GAN \citep{cont2025tailgan} targets tail-risk functionals via VaR/ES-shaped critics. 
CorrGAN \citep{marti2020corrgan} and cCorrGAN \citep{marti2021ccorrgan} generate correlation matrices directly in the elliptope. 
Our approach differs from all of these by demonstrating that the \emph{residual learning} component of financial GANs---typically considered essential---actually hurts OOS correlation fidelity. 
We achieve better results by removing it entirely and using fixed residual covariance instead.

\paragraph{Factor models in asset pricing.}
Our architecture builds on the factor model tradition in finance. 
The APT \citep{ross1976arbitrage} posits that returns are driven by a small number of common factors plus idiosyncratic noise. 
Conditional factor models \citep{kelly2019characteristics} allow time-varying loadings and risk premia. 
Stochastic volatility factor models \citep{pitt1999time, lopes2007factor} extend this to time-varying factor and idiosyncratic variances. 
Our generator can be viewed as a parsimonious stochastic volatility factor model where the volatility dynamics are learned adversarially from macro covariates rather than specified parametrically.

\paragraph{GAN training stability.}
Training stability is a well-known challenge for GANs. 
WGAN-GP \citep{gulrajani2017improved} uses gradient penalty for Lipschitz continuity, while spectral normalization \citep{miyato2018spectral} constrains the discriminator's spectral norm. 
BigGAN \citep{brock2019large} uses orthogonal regularization to prevent mode collapse. 
R1 gradient penalty \citep{mescheder2018training} offers an alternative to WGAN-GP. 
We adopt SN-GAN with hinge loss, which constrains the discriminator through spectral normalization alone and trains reliably at $n_{\text{critic}}=1$ in our setting.

\paragraph{Crisis transmission and correlation dynamics.}
The literature on financial contagion informs our architectural choices. 
\citet{forbes2002no} show that apparent crisis contagion largely reflects continuation of existing linkages through common channels, supporting our decision to route macro stress through factor corrections rather than residual networks. 
\citet{longin2001extreme} demonstrate that correlations increase in bear markets specifically, not just in high-volatility periods. 
These findings motivate our macro-conditioned volatility correction, which adjusts idiosyncratic risk based on current macro state.
